% Options for packages loaded elsewhere
\PassOptionsToPackage{unicode}{hyperref}
\PassOptionsToPackage{hyphens}{url}
\PassOptionsToPackage{dvipsnames,svgnames,x11names}{xcolor}
\documentclass[
  11pt,
  a4paper,
]{article}
\usepackage{xcolor}
\usepackage[margin=1in]{geometry}
\usepackage{amsmath,amssymb}
\setcounter{secnumdepth}{-\maxdimen} % remove section numbering
\usepackage{iftex}
\ifPDFTeX
  \usepackage[T1]{fontenc}
  \usepackage[utf8]{inputenc}
  \usepackage{textcomp} % provide euro and other symbols
\else % if luatex or xetex
  \usepackage{unicode-math} % this also loads fontspec
  \defaultfontfeatures{Scale=MatchLowercase}
  \defaultfontfeatures[\rmfamily]{Ligatures=TeX,Scale=1}
\fi
\usepackage{lmodern}
\ifPDFTeX\else
  % xetex/luatex font selection
  \ifLuaTeX
    \usepackage{luaotfload}
    \directlua{luaotfload.add_fallback("mainfontfallback",{
      "NotoSerifBengali:mode=harf"
    })}
  \fi
  \setmainfont[,RawFeature={fallback=mainfontfallback}]{Latin Modern Roman}
  \ifLuaTeX
    \usepackage{luaotfload}
    \directlua{luaotfload.add_fallback("sansfontfallback",{
      "NotoSansBengali:mode=harf"
    })}
  \fi
  \setsansfont[,RawFeature={fallback=sansfontfallback}]{Latin Modern Sans}
  \ifLuaTeX
    \usepackage{luaotfload}
    \directlua{luaotfload.add_fallback("monofontfallback",{
      "NotoSansBengali:mode=harf"
    })}
  \fi
  \setmonofont[,RawFeature={fallback=monofontfallback}]{DejaVu Sans Mono}
\fi
% Use upquote if available, for straight quotes in verbatim environments
\IfFileExists{upquote.sty}{\usepackage{upquote}}{}
\IfFileExists{microtype.sty}{% use microtype if available
  \usepackage[]{microtype}
  \UseMicrotypeSet[protrusion]{basicmath} % disable protrusion for tt fonts
}{}
\makeatletter
\@ifundefined{KOMAClassName}{% if non-KOMA class
  \IfFileExists{parskip.sty}{%
    \usepackage{parskip}
  }{% else
    \setlength{\parindent}{0pt}
    \setlength{\parskip}{6pt plus 2pt minus 1pt}}
}{% if KOMA class
  \KOMAoptions{parskip=half}}
\makeatother
\usepackage{longtable,booktabs,array}
\usepackage{calc} % for calculating minipage widths
% Correct order of tables after \paragraph or \subparagraph
\usepackage{etoolbox}
\makeatletter
\patchcmd\longtable{\par}{\if@noskipsec\mbox{}\fi\par}{}{}
\makeatother
% Allow footnotes in longtable head/foot
\IfFileExists{footnotehyper.sty}{\usepackage{footnotehyper}}{\usepackage{footnote}}
\makesavenoteenv{longtable}
\setlength{\emergencystretch}{3em} % prevent overfull lines
\providecommand{\tightlist}{%
  \setlength{\itemsep}{0pt}\setlength{\parskip}{0pt}}
% paper_preamble.tex — shared LaTeX preamble for the PCMA/SI paper set.
%
% Compiled with lualatex (NOT pdflatex) because several papers contain
% Bengali script: the CLL paper realises one cognitive state into
% English, Bengali and Banglish, so the Bengali output is EVIDENCE and
% cannot be transliterated away without destroying the point being made.
%
% Font selection is NOT done here. pandoc >= 3 has native
% `mainfontfallback` / `monofontfallback` variables that emit the
% luaotfload fallback registration in the correct position relative to
% unicode-math; doing it by hand in this header lands after unicode-math
% and breaks font loading outright. See tools/build_papers.sh.

\usepackage{amsmath}
\usepackage{amssymb}
\usepackage{mathtools}   % \xleftrightarrow etc. used in cll.md §5
\usepackage{booktabs}

% ---------------------------------------------------------------------
% Unicode symbols appearing in PROSE rather than inside $...$.
%
% The papers were written as markdown and read as markdown, so they use
% real characters — "the ω-boost", "α-child", "ΔR coherence" — instead of
% LaTeX macros. Under unicode-math those codepoints are math-mode-only
% and abort the build ("\symbballowed only in math mode"). Mapping each
% to \ensuremath keeps the markdown readable on GitHub AND compilable,
% which is the whole point of keeping one source.
% ---------------------------------------------------------------------
\usepackage{newunicodechar}

% Greek, lower case
\newunicodechar{α}{\ensuremath{\alpha}}
\newunicodechar{β}{\ensuremath{\beta}}
\newunicodechar{γ}{\ensuremath{\gamma}}
\newunicodechar{δ}{\ensuremath{\delta}}
\newunicodechar{ε}{\ensuremath{\varepsilon}}
\newunicodechar{η}{\ensuremath{\eta}}
\newunicodechar{θ}{\ensuremath{\theta}}
\newunicodechar{κ}{\ensuremath{\kappa}}
\newunicodechar{λ}{\ensuremath{\lambda}}
\newunicodechar{μ}{\ensuremath{\mu}}
\newunicodechar{π}{\ensuremath{\pi}}
\newunicodechar{ρ}{\ensuremath{\rho}}
\newunicodechar{σ}{\ensuremath{\sigma}}
\newunicodechar{τ}{\ensuremath{\tau}}
\newunicodechar{φ}{\ensuremath{\varphi}}
\newunicodechar{χ}{\ensuremath{\chi}}
\newunicodechar{ψ}{\ensuremath{\psi}}
\newunicodechar{ω}{\ensuremath{\omega}}
\newunicodechar{ν}{\ensuremath{\nu}}
\newunicodechar{ξ}{\ensuremath{\xi}}
\newunicodechar{ζ}{\ensuremath{\zeta}}
\newunicodechar{ι}{\ensuremath{\iota}}

% Greek, upper case
\newunicodechar{Δ}{\ensuremath{\Delta}}
\newunicodechar{Φ}{\ensuremath{\Phi}}
\newunicodechar{Σ}{\ensuremath{\Sigma}}
\newunicodechar{Ψ}{\ensuremath{\Psi}}
\newunicodechar{Ω}{\ensuremath{\Omega}}
\newunicodechar{Γ}{\ensuremath{\Gamma}}
\newunicodechar{Λ}{\ensuremath{\Lambda}}
\newunicodechar{Θ}{\ensuremath{\Theta}}
\newunicodechar{Π}{\ensuremath{\Pi}}

% Operators and relations
\newunicodechar{×}{\ensuremath{\times}}
\newunicodechar{·}{\ensuremath{\cdot}}
\newunicodechar{±}{\ensuremath{\pm}}
\newunicodechar{−}{\ensuremath{-}}
\newunicodechar{≈}{\ensuremath{\approx}}
\newunicodechar{≤}{\ensuremath{\leq}}
\newunicodechar{≥}{\ensuremath{\geq}}
\newunicodechar{≠}{\ensuremath{\neq}}
\newunicodechar{∈}{\ensuremath{\in}}
\newunicodechar{∂}{\ensuremath{\partial}}
\newunicodechar{∇}{\ensuremath{\nabla}}
\newunicodechar{√}{\ensuremath{\sqrt{}}}
\newunicodechar{∞}{\ensuremath{\infty}}
\newunicodechar{∝}{\ensuremath{\propto}}
\newunicodechar{∑}{\ensuremath{\sum}}
\newunicodechar{∫}{\ensuremath{\int}}
\newunicodechar{ℝ}{\ensuremath{\mathbb{R}}}
\newunicodechar{ℕ}{\ensuremath{\mathbb{N}}}
\newunicodechar{‖}{\ensuremath{\|}}
\newunicodechar{⊕}{\ensuremath{\oplus}}
\newunicodechar{⊗}{\ensuremath{\otimes}}
\newunicodechar{∀}{\ensuremath{\forall}}
\newunicodechar{∃}{\ensuremath{\exists}}
\newunicodechar{¬}{\ensuremath{\neg}}
\newunicodechar{∧}{\ensuremath{\wedge}}
\newunicodechar{∨}{\ensuremath{\vee}}

% Arrows
\newunicodechar{→}{\ensuremath{\rightarrow}}
\newunicodechar{←}{\ensuremath{\leftarrow}}
\newunicodechar{↔}{\ensuremath{\leftrightarrow}}
\newunicodechar{↑}{\ensuremath{\uparrow}}
\newunicodechar{↓}{\ensuremath{\downarrow}}
\newunicodechar{⇒}{\ensuremath{\Rightarrow}}
\newunicodechar{⇐}{\ensuremath{\Leftarrow}}
\newunicodechar{⇔}{\ensuremath{\Leftrightarrow}}

% Super/subscripts written as literal characters
\newunicodechar{²}{\ensuremath{^{2}}}
\newunicodechar{³}{\ensuremath{^{3}}}
\newunicodechar{⁴}{\ensuremath{^{4}}}
\newunicodechar{⁵}{\ensuremath{^{5}}}
\newunicodechar{⁶}{\ensuremath{^{6}}}
\newunicodechar{₀}{\ensuremath{_{0}}}
\newunicodechar{₁}{\ensuremath{_{1}}}
\newunicodechar{₂}{\ensuremath{_{2}}}
\newunicodechar{₃}{\ensuremath{_{3}}}
\newunicodechar{ᵀ}{\ensuremath{^{\mathsf{T}}}}

% Status glyphs. The papers use emoji as status markers in tables
% (implemented / failed / warning). Mapped to typographic equivalents:
% a colour emoji font in lualatex is unreliable, and the MEANING is what
% carries — a reader needs to see "this one failed", not a red cross.
\newunicodechar{✅}{\ensuremath{\checkmark}}
\newunicodechar{✓}{\ensuremath{\checkmark}}
\newunicodechar{❌}{\ensuremath{\times}}
\newunicodechar{✗}{\ensuremath{\times}}
\newunicodechar{⚠}{\textbf{!}}
\newunicodechar{🧩}{\textbf{[piece]}}
\newunicodechar{📋}{\textbf{[list]}}
\newunicodechar{🟡}{\textbf{[partial]}}
\newunicodechar{🔴}{\textbf{[dead]}}
\newunicodechar{🟢}{\textbf{[live]}}
\newunicodechar{▼}{\ensuremath{\blacktriangledown}}
\newunicodechar{►}{\ensuremath{\blacktriangleright}}
\newunicodechar{□}{\ensuremath{\square}}
\newunicodechar{█}{\rule{0.6em}{0.9ex}}
\newunicodechar{•}{\textbullet}
\newunicodechar{†}{\textdagger}
\newunicodechar{‑}{-}
\newunicodechar{‐}{-}
\newunicodechar{…}{\ldots}
\newunicodechar{।}{\textbar}   % danda: rendered by the Bengali fallback in text
\newunicodechar{⁻}{\ensuremath{^{-}}}

% Zero-width formatting codepoints. These are invisible in an editor but
% are real characters to LaTeX: an emoji written as "⚠️" is TWO
% codepoints (U+26A0 U+FE0F), so declaring only the visible one leaves a
% variation selector with no definition and the build dies on a symbol
% nobody can see. Both are mapped to nothing.
\newunicodechar{^^^^fe0f}{}    % VARIATION SELECTOR-16
\newunicodechar{^^^^200d}{}    % ZERO WIDTH JOINER
\newunicodechar{^^^^200b}{}    % ZERO WIDTH SPACE

\usepackage{longtable}
\usepackage{array}

% Comparison matrices and cost ledgers are wide. Let them shrink to the
% text block rather than overflow the margin.
\usepackage{adjustbox}

% Code blocks: wrap long lines instead of running off the page. Several
% papers quote real C++ with long identifiers.
\usepackage{fvextra}
\fvset{breaklines=true,breakanywhere=true,fontsize=\small}

\setlength{\parskip}{0.5em}
\setlength{\parindent}{0pt}

\usepackage{fancyhdr}
\pagestyle{fancy}
\fancyhf{}
\fancyfoot[C]{\thepage}
\fancyhead[L]{\small\itshape Synthetic Intelligence — research lineage}
\fancyhead[R]{\small\itshape irfan@furylogic.com}
\renewcommand{\headrulewidth}{0.2pt}
\setlength{\headheight}{14pt}
\setlength{\headsep}{18pt}
\usepackage{bookmark}
\IfFileExists{xurl.sty}{\usepackage{xurl}}{} % add URL line breaks if available
\urlstyle{same}
\hypersetup{
  pdftitle={Voltage-Gated Harmonic Topology (VGHT)},
  pdfauthor={Irfan Mahir --- Independent Researcher, Dhaka, Bangladesh --- irfan@furylogic.com},
  colorlinks=true,
  linkcolor={black},
  filecolor={Maroon},
  citecolor={Blue},
  urlcolor={[HTML]{0B5394}},
  pdfcreator={LaTeX via pandoc}}

\title{Voltage-Gated Harmonic Topology (VGHT)}
\author{Irfan Mahir --- Independent Researcher, Dhaka, Bangladesh ---
irfan@furylogic.com}
\date{}

\begin{document}
\maketitle

{
\hypersetup{linkcolor=}
\setcounter{tocdepth}{3}
\tableofcontents
}
\section{Voltage-Gated Harmonic Topology
(VGHT)}\label{voltage-gated-harmonic-topology-vght}

\textbf{Irfan Mahir} \emph{Independent Researcher, Dhaka, Bangladesh}
irfan@furylogic.com

\emph{A three-mechanism physical substrate for domain-orthogonal
cognition with sparse cross-domain analogical bridging and avalanche
suppression, built entirely from scalar field dynamics --- no rules, no
templates, no domain metadata.}

\begin{center}\rule{0.5\linewidth}{0.5pt}\end{center}

\subsection{0. Motivation}\label{motivation}

The Continuous Morphological Dynamics (CMD) substrate --- thermodynamic
habituation, byte-wave sensory injection, phase-locked syntax emission
--- produces coherent, non-templated output but exhibits two coupled
failure modes:

\begin{enumerate}
\def\labelenumi{\arabic{enumi}.}
\tightlist
\item
  \textbf{Domain leakage.} In the shared-oscillator field, a driven
  medical concept can transfer energy into unrelated conceptual clusters
  (e.g., gaming, legal, code) whenever their macro-node phases happen to
  align. The reader observes a medical prompt producing a legal-sounding
  word in the verb slot, or the middle word of a coherent triple
  stepping outside the input's semantic neighborhood.
\item
  \textbf{Polymath rigidity.} The dual failure --- with sufficient
  emission-floor thresholds, the substrate becomes strictly
  domain-locked, unable to draw analogical parallels between
  distantly-related concepts even when the graph explicitly supports
  them through L6-A analogical bridges.
\end{enumerate}

The Continuous Morphological Dynamics literature treats these two
failures as separate tuning problems (emission floor, wavefront gate,
cluster-energy depletion). Voltage-Gated Harmonic Topology (VGHT)
resolves both under a single physical postulate: \emph{domains are
orthogonal frequency bands, and bridges between them are voltage-gated
channels that only conduct under simultaneous high-amplitude resonance.}
Domain isolation and polymath bridging cease to be competing objectives;
they become a single dynamical regime with a single parameter,
θ\_critical, the tunneling threshold.

The mechanism is inspired by three well-understood physical analogs:
mode-orthogonal spectral separation in coupled-oscillator networks,
voltage-gated ion channels in biological neurons (in particular the
sigmoidal open probability of the NMDA receptor with respect to membrane
potential), and the GABAergic global inhibitory tone that suppresses
run-away excitation in cortical circuits. None of the three requires an
explicit ``domain tag'' or a rule dispatcher; each is a continuous
scalar update evaluated once per tick.

\begin{center}\rule{0.5\linewidth}{0.5pt}\end{center}

\subsection{1. Architecture}\label{architecture}

Let the cavity population be indexed by \texttt{i\ ∈\ {[}0,\ N)},
partitioned into macro-clusters \texttt{C\_1,\ C\_2,\ ...} by the
connected components of the tunnel graph (union-find at bind time; the
same clustering used by HNCA). Let a slot layer of \texttt{S} phase
oscillators sit above the macro layer as in CMD phase-locked emission.

\subsubsection{1.1 Domain Sub-Bands}\label{domain-sub-bands}

Each macro-cluster \texttt{C\_k} receives a \emph{base natural
frequency} drawn from a discrete band-bank:

\begin{verbatim}
ω_k = ω_base(k mod B) + jitter(k)
\end{verbatim}

where \texttt{B} is the number of orthogonal bands (default
\texttt{B\ =\ 6}) and \texttt{ω\_base} is a fixed evenly-spaced ladder
in the range \texttt{{[}ω\_lo,\ ω\_hi{]}},
e.g.~\texttt{{[}8,\ 32{]}\ rad/s} giving band centres at
\texttt{\{8,\ 12,\ 16,\ 20,\ 24,\ 28\}}. The jitter is a small in-band
perturbation (±0.5 rad/s) drawn deterministically from the cluster id,
ensuring reproducibility across seeds while still breaking degeneracy
inside a band.

The verb-material ω-boost from prior CMD work becomes an \emph{in-band
shift} rather than a cross-band jump:

\begin{verbatim}
ω_k ← (1 − p_k)·ω_k + p_k·ω_verb_offset(k)
\end{verbatim}

where \texttt{p\_k\ ∈\ {[}0,\ 1{]}} is the verb-material share of
cluster \texttt{k}'s internal edges
(\texttt{verb\_pull\_for(rel)\ /\ \textbar{}internal\ edges\textbar{}}),
and \texttt{ω\_verb\_offset(k)} is the \emph{upper edge} of cluster
\texttt{k}'s home band --- so verb-heavy clusters rise to the top of
their band without colliding with a neighbouring band. This preserves
both the domain orthogonality of Sub-Bands and the verb-slot phase-lock
property from the prior implementation.

\textbf{Consequence.} Two clusters in different bands share almost no
common Kuramoto phase drift, because their \texttt{ω} differs by at
least a band-width and the coupling terms
\texttt{sin(θ\_j\ −\ θ\_i)·A\_j} average to zero over any window longer
than the beat period. Domain interference during ordinary conversation
is suppressed by frequency orthogonality alone --- no relation table, no
cluster tag, no rule.

\subsubsection{1.2 Voltage-Gated L6-A
Tunnels}\label{voltage-gated-l6-a-tunnels}

L6-A tunnels are the analogical bridges discovered by structural
resonance between cross-cluster pairs (Phase 1.3 online-L6A-discovery of
the CMD line of work). In VGHT, every \emph{inter-cluster} tunnel ---
and every tunnel discovered by L6-A regardless of cluster membership ---
is marked \emph{voltage-gated}. Its force is scaled at every tick by a
sigmoidal conductance:

\begin{verbatim}
C(A_a, A_b) = C_max · σ(κ · (A_a · A_b − θ_critical))
\end{verbatim}

where \texttt{A\_a\ =\ \textbar{}x\_a\textbar{}} and
\texttt{A\_b\ =\ \textbar{}x\_b\textbar{}} are the instantaneous
displacements of the two endpoint cavities, \texttt{σ} is the logistic
function, \texttt{κ} is a steepness parameter (default \texttt{12}), and
\texttt{θ\_critical} is the tunneling threshold (default \texttt{0.05}
in the field's natural units). \texttt{C\_max\ =\ 1}.

The force delivered to endpoints is scaled elementwise by \texttt{C}:

\begin{verbatim}
F_a ← F_a · C(A_a, A_b)
F_b ← F_b · C(A_a, A_b)
\end{verbatim}

\textbf{Consequence.} When either endpoint is quiescent,
\texttt{A\_a\ ·\ A\_b\ ≈\ 0}, \texttt{C\ ≈\ σ(−κθ\_critical)\ ≈\ 0}, and
the bridge presents effectively infinite resistance. The two clusters
exchange no energy. When both endpoints are simultaneously excited to a
level whose product exceeds \texttt{θ\_critical}, \texttt{C\ →\ 1}
sharply and the bridge conducts fully. This is the
resonant/quantum-tunneling analog explicitly identified in the
motivating architecture: the door opens only when the room on both sides
is loud enough to make the door irrelevant.

Because both endpoints must resonate at once, spurious cross-domain
leakage --- a stray phase alignment on one side of a bridge --- cannot
open the channel. Only genuine dual excitation, produced by an input
whose structural resonance actually spans two domains, does so. This is
the polymath spark.

\subsubsection{1.3 Global Viscosity (GABA
analog)}\label{global-viscosity-gaba-analog}

The system tracks its total kinetic energy each tick:

\begin{verbatim}
E_total = ½ Σ_i m_i · v_i²
\end{verbatim}

An effective friction coefficient is derived as a saturating function of
\texttt{E\_total}:

\begin{verbatim}
μ_eff = μ_0 + k_avalanche · (1 − exp(−E_total / E_scale))
\end{verbatim}

with defaults \texttt{k\_avalanche\ =\ 0.10}, \texttt{E\_scale\ =\ 5.0}.
\texttt{μ\_eff} replaces the constant friction in the cavity
velocity-decay step and in the tunnel-damping term of that tick.

\textbf{Consequence.} Under normal, sparse activation \texttt{E\_total}
is small, \texttt{μ\_eff\ ≈\ μ\_0}, and the field evolves as before.
Under an avalanche --- many nodes firing simultaneously through spurious
phase alignment --- \texttt{E\_total} rises rapidly, \texttt{μ\_eff}
saturates toward \texttt{μ\_0\ +\ k\_avalanche}, and the excess kinetic
energy is dissipated in the same tick. Only the field components whose
amplitudes are large enough to overcome the elevated friction survive;
the noise floor is damped out. This is the physical analog of GABAergic
inhibitory tone.

Unlike the raw \texttt{exp(E\_total)} formulation, the bounded
\texttt{1\ −\ exp(−E\_total\ /\ E\_scale)} form cannot diverge,
guaranteeing numerical stability under any input, however extreme.

\begin{center}\rule{0.5\linewidth}{0.5pt}\end{center}

\subsection{2. Interaction of the Three
Mechanisms}\label{interaction-of-the-three-mechanisms}

The three mechanisms are not independent tuning knobs; they are coupled
by design.

\begin{itemize}
\tightlist
\item
  \textbf{Sub-Bands} determine which cavities \emph{can} phase-lock
  together in the absence of external excitation. Same-band cavities
  freely couple; cross-band cavities decouple by frequency alone.
\item
  \textbf{Voltage-gated bridges} re-establish cross-band coupling ---
  but only when both endpoints are already excited to a level that
  represents genuine resonance with the input. Without such excitation,
  the bridge is inert regardless of the graph's connectivity.
\item
  \textbf{Global viscosity} suppresses false excitation. When a spurious
  cascade begins to raise \texttt{E\_total}, μ\_eff rises and the
  cascade dies before it can meet the
  \texttt{A\_a\ ·\ A\_b\ \textgreater{}\ θ\_critical} condition needed
  to open any bridge.
\end{itemize}

The joint regime is: \emph{within-band coupling is free and continuous,
cross-band coupling is possible but expensive, and pathological
cross-band cascades are physically forbidden by dissipation.} Domain
isolation and polymath bridging cease to be tradeoffs.

\begin{center}\rule{0.5\linewidth}{0.5pt}\end{center}

\subsection{3. Physical Interpretation}\label{physical-interpretation}

Each of the three mechanisms has a direct biological analog and an
independent literature suggesting the same functional role.

\begin{enumerate}
\def\labelenumi{\arabic{enumi}.}
\tightlist
\item
  \textbf{Frequency-band orthogonality} appears in cortical cognition as
  the observation that different cognitive processes (attention, working
  memory, motor planning) show distinguishable EEG band signatures (α,
  β, γ). Cortical microcircuits couple preferentially within-band, and
  cross-band coupling is a marker of deliberate binding (theta-gamma
  phase-amplitude coupling).
\item
  \textbf{Voltage-gated channels} are the mechanism by which NMDA
  receptors implement coincidence detection: their conductance is
  sigmoidal in the product of presynaptic glutamate release and
  postsynaptic depolarization, mirroring
  \texttt{C(A\_a,\ A\_b)\ =\ σ(A\_a\ ·\ A\_b\ −\ θ)} structurally.
\item
  \textbf{GABAergic inhibition} provides the global tone that keeps
  cortical excitability in the sparse-firing regime. Loss of GABAergic
  inhibition produces exactly the avalanche pattern VGHT's viscosity
  term suppresses: dense synchronous firing across unrelated cortical
  areas, the neural correlate of a seizure.
\end{enumerate}

VGHT does not claim to be a model of any of these biological mechanisms.
It claims that a substrate exhibiting the \emph{dynamical signatures} of
all three simultaneously --- with only three scalar parameters ---
provides the minimum architectural surface area on which
domain-orthogonal cognition with sparse analogical bridging is a stable
regime, rather than an accidental artifact of parameter tuning.

\begin{center}\rule{0.5\linewidth}{0.5pt}\end{center}

\subsection{4. Implementation Contract}\label{implementation-contract}

\begin{itemize}
\tightlist
\item
  \textbf{No if/else branches} in the tick loop that inspect a relation
  label, a domain tag, or a cluster identifier. All three mechanisms are
  pure scalar arithmetic on existing per-cavity or per-tunnel fields.
\item
  \textbf{No new I/O boundary code} for domain metadata. Sub-Bands are
  assigned at bind time from cluster id, computed by the union-find that
  HNCA already runs. Voltage gating is a per-tunnel flag set once at
  bind for cross-cluster tunnels (and for any L6-A tunnel regardless of
  cluster membership). Viscosity is derived each tick from the same
  \texttt{total\_kinetic} accumulator already maintained for the RIFA IC
  pass.
\item
  \textbf{All defaults ON.} VGHT is not a feature flag. The three
  mechanisms are the substrate's normal dynamics.
\end{itemize}

\begin{center}\rule{0.5\linewidth}{0.5pt}\end{center}

\subsection{5. Parameters}\label{parameters}

\begin{longtable}[]{@{}lll@{}}
\toprule\noalign{}
Symbol & Default & Meaning \\
\midrule\noalign{}
\endhead
\bottomrule\noalign{}
\endlastfoot
\texttt{B} & 6 & Number of domain sub-bands \\
\texttt{ω\_lo}, \texttt{ω\_hi} & 8, 32 rad/s & Band ladder range \\
\texttt{κ} & 12 & Voltage-gate steepness \\
\texttt{θ\_critical} & 0.05 & Voltage-gate tunneling threshold \\
\texttt{k\_avalanche} & 0.10 & Peak added friction under avalanche \\
\texttt{E\_scale} & 5.0 & Kinetic-energy saturation scale \\
\end{longtable}

These are the only knobs. All other quantities are derived.

\begin{center}\rule{0.5\linewidth}{0.5pt}\end{center}

\subsection{6. Expected Empirical
Signature}\label{expected-empirical-signature}

A VGHT-enabled substrate is expected to show:

\begin{itemize}
\tightlist
\item
  \textbf{Single-domain prompts} produce single-domain phase-lock
  emissions with no cross-domain leakage into any slot, even after many
  turns of a single-domain conversation. The wavefront gate and the
  sub-band orthogonality are jointly responsible.
\item
  \textbf{Cross-domain prompts} --- where the input drives cavities in
  two different bands to simultaneous high amplitude --- open the
  voltage-gated bridges between those bands within a few ticks,
  producing a genuine polymath emission that spans domains.
\item
  \textbf{Injection of a single garbage or nonsense token} produces at
  worst its own auto-spawned label in output; the avalanche suppressor
  prevents its stray phase drift from recruiting unrelated cavities.
\item
  \textbf{Latency to cross-domain emission} is measurable: it should
  correspond to the time required for the voltage-gate sigmoid to
  saturate under paired excitation, i.e., a few tens of ticks. This
  latency is the phenomenological signature of the polymath spark and is
  measurable with the existing readout infrastructure.
\end{itemize}

\begin{center}\rule{0.5\linewidth}{0.5pt}\end{center}

\subsection{7. Relation to Prior CMD
Work}\label{relation-to-prior-cmd-work}

VGHT subsumes and refines three earlier CMD components without
discarding them.

\begin{itemize}
\tightlist
\item
  Phase 2.2 rubber-tunnel band-pass filters are a special case of
  voltage-gated tunnels with \texttt{θ\_critical\ =\ 0} (always
  conducting) and a Gaussian gate on \texttt{ω} alone. VGHT's sigmoidal
  amplitude gate is strictly more expressive.
\item
  Phase 3.1 multi-ontology substrates rely on manually loaded disjoint
  axiom files to achieve domain separation. VGHT achieves the same
  isolation from a single ontology by frequency orthogonality alone, and
  further allows analogical bridging without touching the ontology.
\item
  The cluster-energy depletion term in phase-locked emission remains a
  fallback for the readout stage. When VGHT is fully active, its
  contribution to output diversity is redundant with the sub-band
  separation.
\end{itemize}

\begin{center}\rule{0.5\linewidth}{0.5pt}\end{center}

\subsection{8. Failure Modes}\label{failure-modes}

The design has three known failure modes worth documenting explicitly.

\begin{enumerate}
\def\labelenumi{\arabic{enumi}.}
\tightlist
\item
  If the cluster count is less than \texttt{B}, some bands go unused,
  which is harmless but wastes spectral capacity.
\item
  If \texttt{θ\_critical} is set too high, cross-domain bridges never
  open even for well-posed polymath prompts. If set too low, the system
  reduces to the pre-VGHT shared-oscillator field.
\item
  If \texttt{k\_avalanche} is set too high, meaningful large-amplitude
  emissions are damped alongside noise; the system becomes
  under-expressive. Empirical tuning against a held-out prompt set is
  suggested.
\end{enumerate}

\begin{center}\rule{0.5\linewidth}{0.5pt}\end{center}

\subsection{9. Summary}\label{summary}

VGHT closes the domain-isolation / polymath-bridging tradeoff by mapping
domains to orthogonal frequency bands, making cross-domain bridges
voltage-gated by joint amplitude resonance, and adding a global
energy-dependent viscosity that dissipates avalanches before they can
spuriously open bridges. All three mechanisms are scalar physics with a
total of six parameters. The tick loop remains free of if/else on
relations, domain tags, or cluster identifiers. The result is a
substrate that is simultaneously domain-locked in default operation and
polymath-capable when its input carries structural resonance across
bands --- the physical realization of the intuition that a polymath is
not a system that mixes domains randomly, but one whose domains are
ordinarily silent to each other and speak only when the same underlying
pattern excites both.

\end{document}
